\section{Route-A interpretation and claim boundary}

The exact incidence result strengthens obstruction OBR-011.  The obstruction is
not that the factorized operator fails to exist: each elementary factor keeps
its declared positive-real normalization and factorized Plancherel unitarity.
The obstruction is narrower.  A single globally nondegenerate type-I phase,
obtained by eliminating \((q_1,q_2)\) in one chart, cannot represent the whole
stationary relation.  A later phase analysis must introduce explicit charts,
transition amplitudes, and a signed Maslov ledger.

The Route-A evaluation is
\[
\text{analytic tuple}=(\mathrm{A1\_WEAK},\mathrm{A2\_FAIL},
\mathrm{A3\_FAIL},\mathrm{A4\_NATURAL\_QUANTIZATION}),
\]
while the Riemann-target tuple remains
\[
 (\mathrm{A1\_WEAK},\mathrm{A2\_FAIL},\mathrm{A3\_FAIL},\mathrm{A4\_FAIL}).
\]
The scoped verdict is \texttt{GO\_WITH\_LIMITATIONS}, with the audit then stopped.

No determinant, spectrum, periodic-orbit census, arithmetic weight law,
completed-\(\xi\) divisor, self-adjoint realization, or zero computation follows
from this certificate.  In particular, a real classical multiplier is not a
Maslov assignment, and the local cubic witness is not evidence for a spectral
zero.  Route B is therefore not authorized.
